\section{Construct a Binary Tree from Preorder and Inorder}
\label{sec:trees-construct-pre-in}

\problemstatement{Given the preorder and inorder traversals of a binary
tree with unique values, reconstruct and return the tree.}

\begin{patternbox}
\emph{``Rebuild a tree from two traversals''} is a divide-and-conquer:
\textbf{preorder's first element is always the root}, and locating that
root inside the inorder array \textbf{splits inorder into the left subtree
and the right subtree}. Recurse on the two halves. A hash map from value to
inorder index turns the ``locate the root'' step from $O(n)$ into $O(1)$.
\end{patternbox}

\subsection*{Understanding the problem carefully}

Preorder visits root, then all of left, then all of right -- so its first
value is the root. Inorder visits left, root, right -- so once we know the
root's value, everything to its left in the inorder array is the left
subtree and everything to its right is the right subtree. The count of
left-subtree nodes (from inorder) tells us exactly how many of the
following preorder elements belong to the left subtree, letting us slice
preorder too. Recursing reconstructs both subtrees.

\begin{ideabox}[Preorder Gives the Root; Inorder Gives the Split]
Take the next preorder value as the current root. Find it in inorder at
position $p$: the $p$ elements left of it form the left subtree, the rest
form the right subtree. Consume preorder left-to-right (root, then build
left, then build right). An index map makes finding $p$ instant.
\end{ideabox}

\subsection*{Algorithm}

\begin{enumerate}
  \item Build \textit{pos}: value $\to$ its index in inorder.
  \item \textsc{Build}(\textit{inLo}, \textit{inHi}) using a shared
        preorder cursor \textit{preIdx} $=0$:
  \begin{enumerate}
    \item If \textit{inLo} $>$ \textit{inHi}, return null.
    \item \textit{rootVal} $=$ preorder[\textit{preIdx}++]; make a node.
    \item $p=\textit{pos}[\textit{rootVal}]$.
    \item node.left $=$ \textsc{Build}(\textit{inLo}, $p-1$); node.right
          $=$ \textsc{Build}($p+1$, \textit{inHi}).
    \item Return node.
  \end{enumerate}
\end{enumerate}

\subsection*{Dry run}

\dryrun{preorder $=[3,9,20,15,7]$, inorder $=[9,3,15,20,7]$.}

\begin{itemize}
  \item Root $=3$ (preorder[0]); in inorder at index $1$: left $=[9]$,
        right $=[15,20,7]$.
  \item Left: root $=9$; inorder slice $[9]$ has no children -- leaf.
  \item Right: root $=20$ (next preorder); in inorder index $3$: left
        $=[15]$, right $=[7]$; both leaves.
  \item Reconstructed: $3$ with left $9$ and right $20(15,7)$.
\end{itemize}

\subsection*{C++ implementation}

\begin{lstlisting}
#include <bits/stdc++.h>
using namespace std;

struct TreeNode {
    int val;
    TreeNode* left;
    TreeNode* right;
    TreeNode(int v) : val(v), left(nullptr), right(nullptr) {}
};

int preIdx;

TreeNode* build(vector<int>& preorder, int inLo, int inHi,
                unordered_map<int,int>& pos) {
    if (inLo > inHi) return nullptr;
    int rootVal = preorder[preIdx++];
    TreeNode* node = new TreeNode(rootVal);
    int p = pos[rootVal];
    node->left  = build(preorder, inLo, p - 1, pos);
    node->right = build(preorder, p + 1, inHi, pos);
    return node;
}

TreeNode* buildTree(vector<int>& preorder, vector<int>& inorder) {
    unordered_map<int,int> pos;
    for (int i = 0; i < (int)inorder.size(); i++) pos[inorder[i]] = i;
    preIdx = 0;
    return build(preorder, 0, inorder.size() - 1, pos);
}

void inorderPrint(TreeNode* n) {
    if (!n) return;
    inorderPrint(n->left);
    cout << n->val << " ";
    inorderPrint(n->right);
}

int main() {
    vector<int> preorder = {3, 9, 20, 15, 7};
    vector<int> inorder  = {9, 3, 15, 20, 7};
    TreeNode* root = buildTree(preorder, inorder);
    inorderPrint(root);            // should reprint the inorder
    cout << "\n";
    return 0;
}
\end{lstlisting}

\begin{complexitybox}
\textbf{Time Complexity.} $O(n)$ -- each node is created once and the
inorder position is an $O(1)$ map lookup. \textbf{Space Complexity.}
$O(n)$ for the map, plus $O(h)$ recursion.
\end{complexitybox}

\begin{takeawaybox}
Reconstruction hinges on two facts: one traversal names the root, another
splits the rest into left and right. The shared preorder cursor consumed
left-to-right is the clean way to avoid slicing arrays at every recursive
call.
\end{takeawaybox}
